\chapter{Programacion lineal}

\section{Modelo canonico}
Un problema de programacion lineal (PL) optimiza una funcion lineal sobre una
region factible descrita por igualdades o desigualdades lineales
\cite{dantzig1963linear,bertsimas1997introduction}.

\begin{definicion}
Sea \(x \in \R^n\) el vector de variables de decision, \(c \in \R^n\), \(A \in
\R^{m \times n}\) y \(b \in \R^m\). Una forma estandar de maximizacion es
\[
  \max \; c^\top x
  \quad \text{sujeto a} \quad
  Ax \leq b,\; x \geq 0.
\]
\end{definicion}

\section{Ejemplo QSB: mezcla de produccion}
Considere dos productos \(X_1\) y \(X_2\). Cada unidad de \(X_1\) aporta 50
unidades monetarias y cada unidad de \(X_2\) aporta 60. Dos recursos limitan la
produccion:
\[
\begin{aligned}
  \max \quad & 50X_1 + 60X_2 \\
  \text{s.a.}\quad & 2X_1 + 3X_2 \leq 180,\\
  & 3X_1 + 2X_2 \leq 150,\\
  & X_1, X_2 \geq 0.
\end{aligned}
\]

El archivo normalizado usado por \qsb{} esta en
\texttt{manual/examples/lp\_sample.json}. La solucion fue obtenida con:

\begin{lstlisting}
swift run qsb solve-json manual/examples/lp_sample.json
\end{lstlisting}

\begin{table}[h]
\centering
\caption{Solucion del ejemplo de programacion lineal}
\begin{tabular}{lr}
\toprule
Cantidad & Valor \\
\midrule
\(X_1\) & 18 \\
\(X_2\) & 48 \\
Valor objetivo & 3780 \\
\bottomrule
\end{tabular}
\end{table}

\section{Interpretacion}
La solucion usa completamente ambos recursos:
\[
2(18)+3(48)=180,\qquad 3(18)+2(48)=150.
\]
Por tanto, las dos restricciones son activas. En una version ampliada del
manual, este ejemplo deberia complementarse con precios sombra, rangos de
sensibilidad y una explicacion geometrica de los vertices factibles.
